\section*{Chapter \ref{ch:iv2:credit-and-xva} --- Credit and XVA}
\markboth{\protect\omTocMark{}{Answers: Credit and XVA}}{\protect\omTocMark{}{Answers: Credit and XVA}}

\begin{solution}{iq:iv2:credit-and-xva:1}
The hazard rate is about 3.33\% a year, and the chance of default within five years is about 15.4\%.

This is the \index{credit triangle}credit triangle: spread = hazard $\times$ (1 $-$ recovery), so $h = 0.02/0.6$. Survival is $e^{-5h} = e^{-1/6} = 0.846$.

The logic is a fair exchange. Each year, the premium is the spread. The expected payout is the chance of default times the loss given default.

\iqwrong{Multiplying $5 \times 3.33\% = 16.7\%$ and calling it the probability. That is close here, but it ignores survival and can go past 100\% for a name in distress.}
\end{solution}

\begin{solution}{iq:iv2:credit-and-xva:2}
It is 3.92\% in the first year, 3.77\% in the second as seen from today, and 3.92\% again given survival of the first.

$1 - e^{-0.04} = 0.0392$. From today, the second year needs survival and then default: $e^{-0.04} - e^{-0.08} = 0.0377$. A constant hazard has no memory, so the conditional figure repeats.

\iqwrong{Saying the second year must be riskier because time has passed. With a constant hazard it is not. The unconditional number is lower only because the company may already be gone.}
\end{solution}

\begin{solution}{iq:iv2:credit-and-xva:3}
With no recovery it is worth 0.7788, a spread of exactly 3\%. With 40\% recovery it is worth 0.8292, a spread of about 1.75\%.

No recovery: $e^{-(0.02 + 0.03) \times 5} = e^{-0.25}$. The spread is the hazard rate.

With recovery of face at maturity: $e^{-0.1}\,(e^{-0.15} + 0.4(1 - e^{-0.15})) = 0.9048 \times 0.9164$. The spread is $-\ln(0.9164)/5 = 1.746\%$, a little below $0.6 \times 3\% = 1.8\%$.

\iqwrong{Using 1.8\% as exact. The triangle is the running-premium result. For a five-year zero with recovery at maturity, it is only an approximation, and it gets worse with higher hazard and longer life.}
\end{solution}

\begin{solution}{iq:iv2:credit-and-xva:4}
You pay \$50,000 a quarter, and on default you receive \$7 million. You are short credit: you gain when the name gets worse.

The annual premium is $0.02 \times \$10$ million $= \$200{,}000$. The payout is $(1 - 0.3) \times \$10$ million. Accrued premium up to the default date is settled as well.

\iqwrong{Calling a protection buyer ``long'', as if it were an option holder. On a credit desk, long credit means carrying the risk, like owning the bond. So the protection seller is long.}
\end{solution}

\begin{solution}{iq:iv2:credit-and-xva:5}
The protection buyer pays the seller about 4.3 points up front. That is about \$430,000 on \$10 million.

Since 2009, North American contracts have traded with fixed coupons of 100 or 500 basis points, and the difference is settled up front. The buyer pays 100 a year for protection worth 200, so the gap is paid at the start.

The hazard is $h = 0.02/0.6 = 3.33\%$. The risky annuity is $(1 - e^{-(r + h)T})/(r + h)$ with $r + h = 6.33\%$, which is 4.29. So the up-front payment is $(0.02 - 0.01) \times 4.29 = 4.29\%$.

\iqwrong{Multiplying by 5 years to get 5 points. The extra premium stops at default and is discounted, so the factor is the risky annuity, not the life.}
\end{solution}

\begin{solution}{iq:iv2:credit-and-xva:6}
At 200, the up-front payment is 4.29 points at 40\% recovery and 4.37 at 20\%. At 2,000, it is 43.8 points and 51.1. So recovery moves the first by 0.1 point and the second by more than 7.

The formula is the same: up-front = (spread $-$ coupon) $\times$ risky annuity, with $h = s/(1 - R)$.

For a name at 200, the hazard moves from 3.33\% to 2.5\%, and the annuity moves only from 4.29 to 4.37. And the spread gap is small, so it hardly matters.

For a name at 2,000, the hazard moves from 33.3\% to 25\%. That moves the annuity from 2.30 to 2.69, and the gap being multiplied is 19 points.

\iqwrong{Saying recovery does not matter because the quote is already in spread. That is true for a par contract, where the up-front is zero whatever the recovery. It fails as soon as the spread is far from the coupon.}
\iqfollow{How would a distressed name be quoted?}{Often in points up front, and sometimes as a price, which leaves less to argue about than a spread that depends on a recovery assumption.}
\end{solution}

\begin{solution}{iq:iv2:credit-and-xva:7}
Buy the bond and buy protection. That is a negative-basis trade, and it earns 40 basis points a year, or \$40,000 on \$10 million, for as long as it lasts.

\textbf{Why the spreads must match.} Hold a par floater with spread $s_b$ and protection at $s_c$, funded at the funding rate. If there is no default, you collect $s_b - s_c$. On a default, you are paid par less recovery by the protection, you deliver or sell the bond at recovery, and you repay par. So you hold no credit risk, and the running gain must be zero if the funding rate is the true financing rate: $s_b = s_c$. With hazard $h$ and recovery of par, both equal $h(1 - R)$.

\textbf{What can go wrong.}
\begin{itemize}
\item Your funding gets dearer, or the repo is pulled. The trade has no credit risk, but it has financing risk, and basis trades get cut when funding tightens.
\item The bond is not par at default, so the payout does not match. A bond above par loses more than protection pays.
\item The contract and the bond can differ on what counts as a default or on what can be delivered.
\item The trade can run at a loss before it converges, and you can be forced out then.
\end{itemize}

\iqwrong{Calling this free money. The extra 40 basis points is what the market pays you for tying up balance sheet and carrying funding risk.}
\end{solution}

\begin{solution}{iq:iv2:credit-and-xva:8}
Both default with probability 1.675\%, and at least one with probability 8.325\%. With probabilities of 5\% and 1\%, the correlation cannot go above about 0.438.

For default indicators, the covariance is $P(\text{both}) - pq$, and the correlation divides that by $\sqrt{p(1 - p)q(1 - q)}$. So
\[
P(\text{both}) = 0.05^2 + 0.3 \times 0.05 \times 0.95 = 0.01675,
\]
and at least one is $0.10 - 0.01675 = 0.08325$.

With unequal probabilities, the joint default cannot go above the smaller one, 1\%. So the highest possible correlation is
\[
\frac{0.01 - 0.0005}{\sqrt{0.05 \times 0.95 \times 0.01 \times 0.99}} = 0.438.
\]

\iqwrong{Assuming a correlation of 1 is always possible. It needs the two events to be the same event, and with different probabilities they cannot be. This also means that ``default correlation'' numbers cannot be compared across names of very different quality.}
\iqfollow{What is the lowest possible correlation here?}{With no overlap, the joint is 0, so the correlation is $-0.0005/0.0217 = -0.023$. Rare events cannot be strongly negatively correlated either.}
\end{solution}

\begin{solution}{iq:iv2:credit-and-xva:9}
Independent: about 2.94\% a year, close to the sum of the five spreads of 0.6\% each. Perfectly correlated: 0.6\%, the same as one name. A second-to-default goes the other way, and rises with correlation.

The chance of at least one default is $1 - 0.99^5 = 4.90\%$, and with recovery of 40\% the expected loss is $0.6 \times 4.90\% = 2.94\%$. With perfect correlation they all go together with probability 1\%, so the loss is $0.6\%$.

The second-to-default needs two defaults. Independent, that is $1 - 0.99^5 - 5 \times 0.01 \times 0.99^4 = 0.098\%$, which is rare. As correlation rises, defaults come in clusters. Under a Gaussian copula, at a correlation of 0.6, the chance of a first default falls from 4.90\% to about 3.7\%, a spread of about 2.2\%, and the chance of a second rises from 0.098\% to about 0.9\%, a spread of about 0.55\%.

\iqwrong{Thinking both basket trades rise with correlation. The first-to-default holder wants defaults spread out, as that makes at least one more likely. The second-to-default holder wants them together.}
\end{solution}

\begin{solution}{iq:iv2:credit-and-xva:10}
$\mathrm{EE}(t) = 0.399 \times 0.01 \times 100 \sqrt{t}\,(5 - t)$. It peaks at $t = 5/3$ years, one third of the life, at about 1.72 on 100 of notional.

At time $t$, the rate has moved by a normal with standard deviation $\sigma\sqrt{t}$, and the swap's value is that move times the remaining annuity, $100(5 - t)$. The expected positive part of a centred normal is $\sigma_X/\sqrt{2\pi} = 0.399\,\sigma_X$.

Setting the derivative of $\sqrt{t}\,(5 - t)$ to zero gives $(5 - t)/(2\sqrt{t}) = \sqrt{t}$, so $t = 5/3$. The value is $0.399 \times 1 \times 1.291 \times 3.333 = 1.72$.

There are two forces. The rate has more time to move, which grows like $\sqrt{t}$. And fewer cash flows are left, which shrinks linearly.

\iqwrong{Saying the peak is at the middle of the life. For a square root against a straight line, it comes at one third.}
\iqfollow{What about a cross-currency swap?}{Its notional is exchanged at the end, so the exposure keeps growing with $\sqrt{t}$ right to maturity and has no such peak.}
\end{solution}

\begin{solution}{iq:iv2:credit-and-xva:11}
The CVA is about 0.0714 on 100 of notional, which is about 7 cents per 100.

$\mathrm{CVA} = (1 - R)\int_0^5 \mathrm{EE}(t)\, h\,\mathrm{d}t$, and the integral of the profile is
\[
0.399\left(5 \cdot \tfrac{2}{3}\, 5^{3/2} - \tfrac{2}{5}\, 5^{5/2}\right) = 5.947.
\]
So $\mathrm{CVA} = 0.6 \times 0.02 \times 5.947 = 0.0714$.

Since $(1 - R)h$ is the spread $s = 1.2\%$, and the average exposure is $5.947/5 = 1.189$, this is $s \times \overline{\mathrm{EE}} \times T = 0.012 \times 1.189 \times 5$. That is the rule you should say aloud: you are paying a CDS premium on your average exposure.

\iqwrong{Using the peak exposure of 1.72 in place of the average. That is about 45\% too high here. Peak exposure is a limit number and not a price.}
\iqfollow{What would discounting at 3\% and survival weighting do?}{Bring it down to about 0.064, as later exposure counts for less.}
\end{solution}

\begin{solution}{iq:iv2:credit-and-xva:12}
With no netting it is 0.798. With netting it is $0.399\sqrt{2 + 2\rho}$: 0, 0.399, 0.564, 0.691 and 0.798 for those five correlations.

With no netting, each trade contributes $0.399$ on its own. With netting, the exposure is on the sum, which is normal with variance $2 + 2\rho$. The expected positive part is 0.399 times its standard deviation.

Netting never makes exposure worse, because $\max(a + b, 0) \le \max(a, 0) + \max(b, 0)$. At a correlation of 1 it gives no gain. At $-1$ the two trades cancel.

\iqwrong{Saying that uncorrelated trades get no netting benefit. They get a lot: 0.564 against 0.798, about 29\% off. Correlation need not be negative.}
\iqfollow{What makes netting valid?}{A legally enforceable master agreement with that counterparty. Without it, a liquidator could claim what you owe while you stand in line for what you are owed.}
\end{solution}

\begin{solution}{iq:iv2:credit-and-xva:13}
It is about \$3.99 million with no collateral and about \$0.79 million with daily margin. It is not zero because the market moves in the 10 days between the last margin call and the close-out.

With no collateral, $0.399 \times \$10$ million $\times \sqrt{1}$. With margin, only the move during the margin period of risk counts: $0.399 \times 10 \times \sqrt{10/252} = 0.795$.

So the ratio is $\sqrt{10/252} \approx 0.2$. Halving the period saves only 29\%.

Also, the margin period is often longer than one might think: disputes, a failure to post, and then replacing the hedge in a stressed market. And exposure over the period tends to come with stress.

\iqwrong{Saying that collateral removes counterparty risk. It reduces it by a factor of about five here, but not to zero. Initial margin is there to cover that gap.}
\end{solution}

\begin{solution}{iq:iv2:credit-and-xva:14}
Without weighting, CVA is 0.0714 and DVA is 0.0357, so the bilateral adjustment is 0.0357 against you. With first-to-default weighting at $e^{-(h_c + h_o)t}$, they are 0.0670 and 0.0335, so 0.0335.

An at-the-money swap has a symmetric profile, so the negative exposure has the same integral, 5.947. So $\mathrm{CVA} = 0.6 \times 0.02 \times 5.947$, and $\mathrm{DVA} = 0.6 \times 0.01 \times 5.947$. The weighting counts only the first of the two defaults, since after that the contract has stopped.

If your own credit worsens, DVA grows, and you book a profit. A firm in trouble posts a gain for being in trouble.

Desks are wary of this for good reasons:
\begin{itemize}
\item You can only realise it by defaulting, or by selling protection on yourself, which you cannot do.
\item You can hedge it only with a proxy, such as a basket of peers, and that adds risk.
\item Capital rules take it out of regulatory capital, so it does not support the balance sheet.
\end{itemize}

Accounting fair value still includes it, as an exit price would reflect your own credit.

\iqwrong{Dropping DVA as simply fictitious. If you quote a price to a counterparty that is safer than you, it will ask for it, so DVA is real in the price even if it is hard to monetise.}
\iqfollow{What if both hazards are equal and the swap is at the money?}{Then CVA and DVA cancel exactly, and you would quote no credit charge at all, with or without the weighting.}
\end{solution}

\begin{solution}{iq:iv2:credit-and-xva:15}
The funding cost at 100 basis points is about 0.0595 on 100 of notional. When the spread is purely credit, at 60 basis points, the funding benefit is 0.0357, which is exactly your DVA. So counting both double-counts your own credit.

\textbf{FVA.} When the client's swap is in the money to you, the hedge is out of the money, and you post collateral on it. You borrow to fund that collateral, and the client posts nothing. So you pay the funding spread on the positive exposure: $0.01 \times 5.947 = 0.0595$.

When the client's swap is out of the money, you receive collateral on the hedge, and that saves you borrowing. So you gain the spread on the negative exposure, which is the funding benefit.

\textbf{The overlap.} With funding spread $s_o = h_o(1 - R) = 0.6\%$, the benefit is $s_o \times 5.947 = 0.0357$. And $\mathrm{DVA} = (1 - R)h_o \times 5.947$, the same number. It is one thing seen twice: you may never repay that debt, which is why your borrowing costs more and also why your liability is worth less.

So a coherent price takes the cost side from funding and the benefit side from either DVA or funding, not both.

The debate is whether FVA belongs in a price at all, as a pure no-arbitrage argument would say funding should not change value. Banks book it because they pay it. If the funding spread is zero, FVA is zero, as it should be.

\iqwrong{Adding CVA, DVA, FVA cost and FVA benefit without looking for overlap. If spreads are all credit, the sum counts your own default twice.}
\iqfollow{Does the funding cost depend on the counterparty?}{Only through survival: the funding need stops if the client defaults. So it is mostly your own spread that matters, which is why two banks will price the same trade differently.}
\end{solution}

\begin{solution}{iq:iv2:credit-and-xva:16}
The expected exposure given default is 1.99 at $\rho = 0$, 4.25 at $0.3$, 6.22 at $0.5$, and 0.62 at $-0.3$. So $\rho = 0.3$ more than doubles the CVA, and $\rho = 0.5$ triples it.

Exposure is large when $Z$ is low, and with positive $\rho$ default is also more likely when $Z$ is low. So you lose most when the counterparty is most likely to go: that is \index{wrong-way risk}wrong-way risk.

Given default, the factor is shifted. The default threshold for 2\% is $-2.05$, and the mean of a normal below that is $-\varphi(c)/\Phi(c) = -2.42$. So $Z$ has a conditional mean of $-2.42\rho$, and $E[\max(-5Z, 0) \mid \text{default}]$ is an integral:
\[
\frac{1}{0.02}\int_{-\infty}^{0} (-5z)\,\varphi(z)\,\Phi\!\left(\frac{c - \rho z}{\sqrt{1 - \rho^2}}\right)\mathrm{d}z.
\]
At $\rho = 0$ it is $5\varphi(0) = 1.99$, the plain number.

\textbf{Wrong-way examples.} Buying protection on a bank from a bank in the same system. Holding a dollar receivable from a borrower in a weak currency, where the currency falls as it fails. Lending against the counterparty's own shares.

\textbf{Right-way example.} An airline that has bought fuel forward from you. You are exposed when fuel falls, and falling fuel helps the airline. Always check the direction: an oil producer that has bought oil forward from you is wrong way, as falling oil raises your exposure and weakens the producer.

\iqwrong{Multiplying average exposure by default probability and stopping. That assumes independence, and it can be off by a factor of three.}
\end{solution}

\begin{solution}{iq:iv2:credit-and-xva:17}
The table gives the expected loss as a share of each tranche.

\begin{center}
\begin{tabular}{lcccc}
\toprule
$\rho$ & 0--3\% & 3--7\% & 7--15\% & 15--100\% \\
\midrule
0.05 & 39.6\% & 0.33\% & 0.00\% & 0.00\% \\
0.2 & 33.5\% & 4.06\% & 0.40\% & 0.00\% \\
0.5 & 21.3\% & 6.70\% & 2.42\% & 0.12\% \\
\bottomrule
\end{tabular}
\end{center}

Equity loses value as correlation falls and gains as it rises. The senior tranche moves the other way. What stays fixed is the pool's expected loss, $0.02 \times 0.6 = 1.2\%$, since a sum of expectations ignores correlation.

Given the market factor $M$, a large pool's loss is deterministic:
\[
L(M) = (1 - R)\,\Phi\!\left(\frac{\Phi^{-1}(p) - \sqrt{\rho}\,M}{\sqrt{1 - \rho}}\right).
\]
The tranche's loss is $E[\min(\max(L - a, 0), b - a)]/(b - a)$, by quadrature over $M$.

There are two limits. As $\rho \to 0$, the pool loses exactly 1.2\%, so the first tranche has lost $1.2/3 = 40\%$, and the rest nothing. As $\rho \to 1$, all default together with probability 2\% and the loss is 60\%, so every tranche up to 60\% has lost 2\%.

So equity is long correlation, and senior is short correlation. The middle tranche moves one way and then the other: it rises here, but must go back toward 2\% as $\rho \to 1$.

\iqwrong{Thinking there is one correlation for all tranches in the market. A quoted model correlation differs by tranche. That is the correlation smile, and it says the one-factor Gaussian copula does not fit.}
\iqfollow{Where does this model fail in a crisis?}{It has no tail dependence, and each pair's correlation is one number. Li's use of the Gaussian copula made default correlation easy to quote, but it does not say that defaults cluster more in bad times.}
\end{solution}

\begin{solution}{iq:iv2:credit-and-xva:18}
Equity is worth 24.15, debt 75.85, the risk-neutral default probability is 18.7\%, and the spread is about 232 basis points. At three months, the spread is about 71 basis points, and at five years about 210.

Equity is a call on the assets. So $d_1 = 1.1376$, $d_2 = 0.8876$, and $E = 100\,\Phi(d_1) - 80 e^{-0.03}\Phi(d_2) = 24.15$. Debt is what is left: $100 - 24.15 = 75.85$. A second route is debt as a riskless bond less a put: $77.64 - 1.78 = 75.85$.

Default is when assets end below 80, so the risk-neutral probability is $\Phi(-d_2) = 18.7\%$.

The spread is $-\ln(75.85/80) - 0.03 = 2.32\%$.

\textbf{Term structure.} At one month, the spread is about 2 basis points, and as maturity goes to zero, it goes to zero. It peaks near two years at about 248, then falls: 210 at five years, 162 at ten.

The short end is the failure. A company can only reach default by a continuous diffusion, and over a short time it can't get there. So the model gives almost nothing for short-term credit, while real short spreads are not close to zero. Models add jumps, or uncertainty about the asset level or the barrier, to fix it.

Also, asset value and volatility are not observed. Practitioners back them out of equity and equity volatility, using $\sigma_E E = \Phi(d_1)\sigma_A V$.

\iqwrong{Reading 18.7\% as a real-world probability. It uses the risk-neutral drift. With a real-world asset drift above the rate, the true probability is lower.}
\iqfollow{What is distance to default?}{It is $(\ln(V/D) + (\mu - \sigma^2/2)T)/(\sigma\sqrt{T})$, the number of standard deviations the assets are from the debt. It is $d_2$ with the real drift in place of the rate.}
\end{solution}

\begin{solution}{iq:iv2:credit-and-xva:19}
The triangle gives 1.67\%, 2.92\% and 4.58\%. The bootstrap gives 1.67\%, 2.95\% and 4.81\%.

\textbf{Triangle.} Each spread gives an average hazard: $0.01/0.6$, $0.015/0.6$, $0.02/0.6$. Forward pieces follow from total hazard: $(3 \times 2.5\% - 1.67\%)/2 = 2.92\%$ and $(5 \times 3.33\% - 3 \times 2.5\%)/2 = 4.58\%$.

\textbf{Bootstrap.} Solve in order. Take the first hazard so that the one-year contract prices to zero. Then, holding it, take the second so the three-year contract prices to zero, then the third. Premium leg: $\sum_i \tfrac14 D_i S_i$ plus half a quarter of accrued on default. Protection leg: $(1 - R)\sum_i D_i (S_{i-1} - S_i)$.

\textbf{Why higher.} A CDS spread is a weighted average of the forward spreads, and the weights are the risky discount factors $D_iS_i$, not equal. Early periods weigh more, because they are less discounted and more likely to be reached. So the triangle's equal time weights underestimate how high the last piece has to be to lift the average to 200.

\textbf{The inverted curve.} At 2,000 and 800, the first year has a hazard of about 33\%. Total hazard over five years by the triangle is $5 \times 13.3\% = 66.7\%$, so the next four years average about 8.3\%. The market says: this name is in immediate danger, and if it gets through that, it will be much safer. A distressed name's long-dated contract is close to a claim on the same near-term event.

\iqwrong{Fitting one flat hazard to each maturity and calling that the curve. Different flat hazards for different maturities are not consistent with one survival curve, and they cannot price a forward-starting contract.}
\iqfollow{What can go wrong in a bootstrap?}{A steeply inverted curve with a high recovery can force a negative forward hazard, which is an arbitrage in the inputs or a sign that the recovery assumption is too high.}
\end{solution}

\begin{solution}{iq:iv2:credit-and-xva:20}
The CVA is about \$110,800. Its sensitivity to the spread is about \$515 per basis point, so buy about \$1.1 million of protection. After the rate move, CVA is \$182,500. After the spread move it is \$160,600. After both it is \$265,400. So the two together add \$33,100 beyond the sum.

\textbf{Value.} With $h = 0.02/0.6 = 3.33\%$ and $\mathrm{EE}(t) = 0.399\sqrt{t}\,(5 - t)$ per 100,
\[
\mathrm{CVA} = 0.6 \int_0^5 h e^{-ht}\,\mathrm{EE}(t)\,\mathrm{d}t = 0.1108 \text{ per } 100.
\]
Without the survival weighting, it would be $0.02 \times 5.947 = 0.119$.

\textbf{Hedge.} Revalue with the spread at 201: CVA rises by \$515. A \$1 million 5-year CDS gains its risky annuity times one basis point: $(1 - e^{-5h})/h \times 10^{-4} \times \$1$ million $= 4.61 \times \$100 = \$461$. So you need $515/461 = \$1.12$ million of protection, which is roughly 1\% of the swap's notional. It is small because exposure is small against notional.

\textbf{Rate move.} If the rate has moved so that the swap is in your favour by 50 basis points, exposure gets a mean $m = 0.005$ and
\[
\mathrm{EE}(t) = 100(5 - t)\,\sigma_t\,\bigl(\varphi(m/\sigma_t) + (m/\sigma_t)\Phi(m/\sigma_t)\bigr), \quad \sigma_t = \sigma\sqrt{t}.
\]
That gives \$182,500. A 50-basis-point move against the swap would give \$64,100, so the sensitivity is not symmetric, and the delta at the money is about \$1,180 per basis point.

\textbf{Both.} Adding the separate moves gives $110{,}800 + 71{,}700 + 49{,}700 = \$232{,}200$, and the true value is \$265,400. The extra \$33,100 is \index{cross-gamma}cross-gamma between rates and credit. When you are in the money the CS01 is \$855, up from \$515. So your spread hedge needs to be \$1.86 million, and it was short by \$0.74 million for exactly the move that hurts.

So you must rebalance the credit hedge when rates move, and the rates hedge when spreads move. And a desk that hedges each risk separately still pays the cross term. With wrong-way risk, the cross term gets larger, as spread and exposure then move together.

\textbf{No traded CDS.} Hedge with a proxy, such as a sector and rating index or a related name. You remove the systematic spread risk, and you keep the idiosyncratic risk. That is a real residual, and you hold reserves or capital against it. You state plainly that the jump to default is not hedged.

\begin{dated}{2026-11}{Why the CVA charge exists}\label{dat:iv2:credit-and-xva:cva}
The Basel Committee's press release on the capital treatment of bilateral counterparty risk says that in the crisis roughly two-thirds of the losses attributed to counterparty credit risk came from CVA losses, and only about one-third from actual defaults. That is why the capital rules charge the mark-to-market movement of CVA and not only the default.
\end{dated}

\iqwrong{Hedging the CVA only for spread, as one number. CVA is an option-like position on the portfolio's mark to market, multiplied by a credit curve. It has rates delta, vega, and cross-gamma, and those are where the losses come from.}
\iqfollow{Does a 20\% rise in rate volatility matter?}{At the money, yes, nearly in proportion: CVA goes to \$133,000, because exposure is linear in volatility. So a CVA desk is short vega in its book's rates: it loses when vol rises.}
\end{solution}
